\documentclass[11pt]{article}
\usepackage[utf8]{inputenc}
\usepackage{amsmath,amssymb}
\usepackage{authblk}
\usepackage[margin=1in]{geometry}
\usepackage[hidelinks]{hyperref}
\usepackage{xurl}
\usepackage{setspace}

\title{RealQM Atomic Spectra\\
\large Radiation as Real Charge Oscillation: Sloshing and Breathing}
\author[1]{Claes Johnson}
\affil[1]{KTH Royal Institute of Technology, SE-100 44 Stockholm, Sweden \\
\texttt{claesjohnson@gmail.com}}
\date{\today}

\begin{document}
\maketitle

\begin{abstract}
Atomic spectra were the starting point of quantum mechanics: the discrete lines a classical atom cannot
explain. Standard quantum mechanics accounts for them \emph{formally} --- the levels are eigenvalues, and a
spectral line is the beat of a \emph{superposition} of two eigenstates, a formal device whose physical
meaning is left open, the same unresolved question as measurement. Real Quantum Mechanics (RealQM) --- which
represents $N$ electrons by $N$ unit charge densities on non-overlapping domains of ordinary
three-dimensional space, meeting at free boundaries --- offers instead a physical account. Its stationary
states are the \emph{stationary points} of a single global energy, at total energies $E_n$; making the wave
functions complex turns radiation into a \emph{real charge-density oscillation} between them. Two mechanisms
arise, and \emph{neither is a superposition}: \emph{sloshing}, a real transition in which the charge moves
between two full states and radiates at $E_n-E_m$; and \emph{breathing}, in which two shells of one state
coexist side by side and beat at $E_j-E_k$. The two forms of helium appear as the two modes --- ortho a
sloshing at $19$~eV (observed $19.8$), para a breathing at $43$~eV. The most visible feature of an atom, carried in the
standard theory by a formal superposition, becomes here a literal oscillation of charge in space; the
quasi-static free boundary and its fast-scale limit are the supporting time-dependent machinery.
\end{abstract}

\setstretch{1.06}

\section{Real-valued RealQM}
\label{sec:static}

RealQM represents $N$ electrons by densities $\psi_j^2$, each carried on its own region
$\Omega_j\subset\mathbb R^3$, the regions non-overlapping, $\Omega_j\cap\Omega_k=\varnothing$
\cite{RealQM}; the nuclei enter as a fixed kernel. The wave functions $\psi_j$ are real, each of unit charge,
\begin{equation}
\int_{\Omega_j}\psi_j^2=1 ,
\label{eq:unit}
\end{equation}
and of finite kinetic energy: each $\psi_j\in H^1(\Omega_j)$, and the total wave function
$\Psi=\sum_j\psi_j\in H^1(\mathbb R^3)$, where $H^1(\Omega)$ denotes the functions on $\Omega$ that are
square-integrable together with their first derivatives. Membership of $\Psi$ in $H^1(\mathbb R^3)$ forces the
$\psi_j$ to \emph{agree on the interfaces between neighbouring domains}; this continuity is the free-boundary
condition (i) below. The ground state minimises the Coulomb energy
\begin{equation}
E=\sum_j\tfrac12\!\int_{\Omega_j}|\nabla\psi_j|^2+\int V_{\mathrm K}\,\rho\,dx
+\sum_{j<k}\iint\frac{\psi_j^2(x)\,\psi_k^2(y)}{|x-y|}\,dx\,dy,
\qquad \rho=\sum_j\psi_j^2 ,
\label{eq:E}
\end{equation}
with $V_{\mathrm K}$ the Coulomb potential of the fixed nuclear kernel and $\rho$ the total electron density.
The three terms are the electron kinetic energy, the electron--kernel attraction, and the electron--electron
repulsion --- the last running over \emph{distinct} domains $j<k$ only. There is \emph{no self-repulsion}: an
electron does not repel its own charge. That confining role is taken over by the kinetic energy
$\tfrac12\int|\nabla\psi_j|^2$, whose pressure keeps each density from collapsing; kinetic pressure balanced
against the kernel attraction and the repulsion of the other electrons fixes the size of each domain. The
energy is minimised over the wave functions \emph{alone}; the partition $\{\Omega_j\}$ --- the free boundary
--- is not a separate variable but an outcome of that minimisation. Minimising each $\psi_j$ on its region,
with its boundary value free, gives on each domain the local (Euler--Lagrange) equation in the total
self-consistent field of the kernel and the \emph{other} electrons, together with its \emph{natural} boundary
condition, the Neumann condition (ii) below; continuity (i) then locates the interface. Each domain carries
its own Poisson problem for its potential, coupled to the others only through the fields they produce.

Two conditions hold at each interface $\Gamma$, with distinct roles.
\begin{itemize}
\item[(i)] \emph{Density continuity.} $\rho_1=\rho_2$ on $\Gamma$; since $\rho_j=\psi_j^2$ and the $\psi_j$
are real, this is $\psi_1=\psi_2$.
\item[(ii)] \emph{Vanishing normal derivative} (homogeneous Neumann, the Bernoulli condition):
$\partial\psi_j/\partial n=0$ on $\Gamma$.
\end{itemize}
Condition (ii) belongs to each domain's \emph{own} problem: the vanishing normal derivative is the
homogeneous, zero-flux boundary condition that closes the Poisson problem on $\Omega_j$, so
that each domain is a self-contained subsystem across which no charge flows and its unit charge is kept.
Condition (i) is the extra free-boundary condition: with each domain solving its own Neumann problem on
whatever region it is given, the interface is \emph{located} where the two densities meet, $\rho_1=\rho_2$.
So (ii) is part of the equation on each domain and (i) fixes the position of the free boundary. With both
imposed the static free-boundary problem is well posed, and it is this problem --- unit charges, density
continuity, and Neumann, coupled by Coulomb --- that reproduces the periodic table,
chemical bonding, and, read for protons and electrons, nuclear binding \cite{RealQM,RealNucleus}. Everything
here is static and real: the density is fixed in time and no charge flows.

\section{Complex-valued RealQM}
\label{sec:complex}

The complex phase is \emph{periodic motion}: a clock $e^{-iEt}$ turning at rate $E$. Its point is radiation
understood as \emph{resonance} --- the agreement of periodic motions. When a charge carries two frequencies
$E_n,E_m$, their beat at $E_n-E_m$ is a real periodic oscillation, and radiation is this oscillation coming
into resonance with the electromagnetic field. A static real density has its phases frozen, no periodic
motion, and cannot resonate; making the wave functions \emph{complex} restores it --- this is what the
complex form is \emph{for}. To reach it, and the other phenomena absent from a static real density (currents,
magnetism), one evolves each domain by the Schr\"odinger equation generated by \eqref{eq:E},
\begin{equation}
i\,\partial_t\psi_j=-\tfrac12\nabla^2\psi_j+\big(V_{\mathrm K}+\phi_j\big)\psi_j ,\qquad
\phi_j(x)=\int\frac{(\rho-\rho_j)(y)}{|x-y|}\,dy ,
\label{eq:tdse}
\end{equation}
with $V_{\mathrm K}+\phi_j$ the total Coulomb potential felt by electron $j$ --- the fixed kernel plus the
Hartree field $\phi_j$ of the \emph{other} electrons (no self-repulsion). This is the \emph{entire}
extension. It costs nothing in the static construction, and for a clean reason: everything the theory couples
through is the density. The free boundary is located by (i)--(ii), which are conditions on
$\rho_j=|\psi_j|^2$; the Coulomb field in \eqref{eq:tdse} is the field of the other electrons and the fixed
kernel. Both are \emph{phase-blind}. Writing
$\psi_j=|\psi_j|\,e^{-i\theta_j}$ changes neither the density, nor the boundary it fixes, nor the energy
\eqref{eq:E}. Complexification does one thing only: it restores the phase $\theta_j$ that the real theory had
frozen to zero. The static real-valued theory of \S\ref{sec:static} is the phase-free section of the same
theory, not a different one.

The phase, once free, is dynamical. A complex $\psi_j$ carries a current
\begin{equation}
\mathbf j_j=\operatorname{Im}(\psi_j^*\nabla\psi_j)=-|\psi_j|^2\,\nabla\theta_j,\qquad
\partial_t\rho_j+\nabla\!\cdot\mathbf j_j=0 ,
\label{eq:current}
\end{equation}
identically zero when $\psi_j$ is real ($\theta_j\equiv0$). The current is what a static real density lacks,
and it is set by the \emph{gradient} of the phase.

\section{Stationary states as stationary points of the energy}
\label{sec:eigen}

The ground state of \S\ref{sec:static} is the \emph{minimum} of the global energy \eqref{eq:E}. In the spirit
of Schr\"odinger, the stationary states of the system are the \emph{stationary points} of that same global
energy --- the ground state the least, the excited states the higher stationary points --- each with a total
energy $E$, the value of \eqref{eq:E} there, that is the state's energy.

Complexified, such a stationary state is carried by clocks
\begin{equation}
\psi_j(x,t)=u_j(x)\,e^{-iE_jt},
\label{eq:stat}
\end{equation}
a stationary solution of the evolution \eqref{eq:tdse}, each domain turning at its own rate $E_j$ with density
$\rho_j=u_j^2$ time-independent and the Coulomb field stationary. There are \emph{no} separate eigenvalue
problems on the individual domains: stationarity of \eqref{eq:E} under variation of $\psi_j$, with the
unit-charge constraint \eqref{eq:unit} carried by a multiplier $E_j$, gives on each domain only the local
condition
\begin{equation}
-\tfrac12\nabla^2 u_j+\big(V_{\mathrm K}+\phi_j\big)u_j=E_j\,u_j\quad\text{on }\Omega_j,\qquad
\frac{\partial u_j}{\partial n}=0\ \text{ on }\partial\Omega_j,
\label{eq:eigen}
\end{equation}
with $V_{\mathrm K}+\phi_j$ the total Coulomb field of the kernel and the other electrons (no self-repulsion),
the Neumann condition (ii) the natural boundary condition of the minimisation, the interfaces located by continuity (i), and
$\rho=\sum_j u_j^2$. The $\psi_j$ are coupled through the shared field and the free boundaries, and $E_j$ is
merely the multiplier of domain $j$'s constraint --- the local face of one global stationarity, not a
domain-by-domain spectral problem.

The \emph{total} energy $E$ --- the value of \eqref{eq:E} --- is distinct from the clock rates: because the
domains share their interactions, $E\neq\sum_jE_j$, and only if all the $E_j$ coincided would the domains
merge into a single global clock $\Psi\,e^{-iEt}$, the textbook stationary state at total energy $E$. What the
ground-state minimisation of \S\ref{sec:static} carries as the Lagrange multiplier of the unit-charge
constraint reappears, in the complex form, as the physical clock rate $E_j$.

Beyond the ground state lie the other stationary points of \eqref{eq:E}, at higher total energy: a discrete
family of \emph{full states} of the system, with total energies $E_n$, $E_0<E_1<\cdots$ --- exactly the
levels Schr\"odinger's stationary problem provides, a one-electron atom its single-electron levels, a
many-electron atom its configurations. Radiation is a real charge oscillation \emph{between} two such full
states, at a total-energy difference $E_n-E_m$ (\S\ref{sec:spectra}). Complexifying thus carries RealQM from
the single energy minimum to the whole family of stationary states.

\section{Two time scales: a quasi-static boundary and per-domain phase clocks}
\label{sec:twoscale}

The eigenstates of \S\ref{sec:eigen} hold the geometry fixed, each domain ticking at its own $E_j$. Letting
the geometry drift --- the domains and their free boundary relaxing \emph{quasi-statically} while each domain
keeps its clock --- gives the general two-scale solution:
\begin{equation}
\psi_j(x,t)=u_j\big(x;\Gamma(t)\big)\,e^{-i\theta_j(t)} ,
\label{eq:ansatz}
\end{equation}
with the geometry $\Gamma(t)$ --- the domains and their free boundary --- moving \emph{quasi-statically},
relaxed at each slow step by the static conditions (i)--(ii) on the instantaneous densities $u_j^2$, while
each domain runs its own \emph{phase clock} $\theta_j(t)=\int E_j\,dt$. The clock is spatially uniform across
a domain, so by \eqref{eq:current} it carries \emph{no current},
\begin{equation}
\nabla\theta_j=0\ \Longrightarrow\ \mathbf j_j=0 ,
\end{equation}
the density $u_j^2$ is stationary at the fast scale, each electron keeps its unit charge automatically, and
the boundary --- which responds only to the density --- has nothing fast to track. The static pair
(i)--(ii) is all it ever sees. There is no fast moving-boundary problem to solve and no equation of motion
to impose on $\Gamma$; the geometry simply follows the slow drift of the densities as in the real theory.

One point makes the phases free. To let different domains tick at different rates, the interface matches the
\emph{density} \eqref{eq:cont}, not the complex value: $|\psi_1|=|\psi_2|$ with the phases unconstrained
across $\Gamma$. Full complex continuity, $\psi_1=\psi_2$, would force $\theta_1=\theta_2$ and hence
$E_1=E_2$, collapsing the clocks to one. Density continuity --- which is what RealQM asks anyway, continuity
of $\psi^2$ being continuity of charge --- liberates them.
\begin{equation}
|\psi_1|=|\psi_2|\ \text{ on }\Gamma,\qquad \partial_n u_j=0\ \text{ on }\Gamma,\qquad
\theta_j\ \text{free across }\Gamma .
\label{eq:cont}
\end{equation}

\section{Coexistence of time scales}
\label{sec:super}

A single uniform clock, alone, has no observable content: it can be gauged away, carries no current, and
leaves every density stationary. The phase becomes physical through the \emph{differences} of energies, and it
does so in two ways --- neither of them a superposition.

\emph{Between two full states} $\Psi_n,\Psi_m$ (\S\ref{sec:eigen}) at total energies $E_n,E_m$, the charge
density physically moves from one to the other --- a real transition --- radiating at $E_n-E_m$ as it goes:
the system \emph{sloshing} between states. Standard quantum mechanics instead writes the atom as a
superposition of global eigenstates,
\begin{equation}
\Psi=a\,\Psi_n\,e^{-iE_nt}+b\,\Psi_m\,e^{-iE_mt},
\label{eq:super}
\end{equation}
a formal object whose physics is unclear or even missing --- what it \emph{is} to be in a superposition of two
energy states is the unresolved question of measurement. RealQM carries no such object: only a real charge
density in motion, whose oscillation at $E_n-E_m$ is the radiation. \emph{Between two domains of a single
state}, an inner and an outer shell coexist side by side, and the relative phase $\theta_j-\theta_k$ advances
at $E_j-E_k$ --- a \emph{breathing}, again no superposition but two definite charges beating against each
other, a beat any coupling that bridges the domains sees. The phase field is thus the register on which the
atom's \emph{time scales} are carried at once: the slow geometric scale of the moving domains and the fast
electronic scales --- the total energies $E_n$ and the domain rates $E_j$ --- and their differences. This is
the content a static real density cannot hold --- spectra, coherence, radiation --- recovered by the single
step of letting $\psi$ be complex.

The reading is standard-QM-friendly. Equation \eqref{eq:ansatz} is the adiabatic (Born--Oppenheimer)
separation of a slow configuration from a fast phase, with $e^{-iE_jt}$ the ordinary stationary-state clock
--- here carried \emph{per electron domain} rather than globally. Within each domain \eqref{eq:tdse} is the
familiar one-body complex Schr\"odinger equation in a self-consistent field, so its whole apparatus ---
split-operator or Crank--Nicolson propagation, resonance and complex-energy methods, time-dependent
mean-field --- transfers unchanged. What is new to RealQM is only the \emph{interface} between domains, and
in the quasi-static reading it is handled by the same static Bernoulli conditions as before.

\section{Two sources of spectra: sloshing and breathing}
\label{sec:spectra}

The spectrum is the most visible aspect of an atom --- how atoms are seen, identified, and were first
discovered --- yet standard quantum mechanics hides the physics of how it arises: the levels drop out of an
eigenvalue problem and the lines from transition probabilities, but what \emph{physically} radiates is never
said. In RealQM radiation is a \emph{real charge-density oscillation}: two coexisting energies make the charge
physically oscillate and radiate, as in Schr\"odinger's reading of the beat of eigenstates, with no collapse
of an amplitude. The complex theory carries two such oscillations, told apart by whether the beat is between
two different \emph{states} of the system or between two \emph{shells of one state}.

\emph{Sloshing --- between two full states.} The system moves between two of its stationary states
(\S\ref{sec:eigen}), at total energies $E_n$ and $E_m$: the charge really sloshes from one configuration to
the other and radiates at $E_n-E_m$. This is the ordinary atomic transition, and it needs no more than a
single electron --- a hydrogen atom sloshes between $1s$ and $2p$, a many-electron atom between its
configurations. It is the clear, primary source of spectra.

\emph{Breathing --- between two shells of one state.} A single multi-shell state can also carry an oscillation
\emph{within} itself: two domains, an inner and an outer shell, coexist side by side at different clock rates
$E_j,E_k$, and their coexistence sustains radiation at $E_j-E_k$, with no transition to another state. This is
a second, more delicate source. Its physics is less clear --- it radiates only if the shells \emph{couple} (a
shared boundary that passes charge, or the Coulomb field), and a spherical inner--outer beat carries no dipole
--- but it is free of the superposition enigma: the two shells are two definite charges side by side, not one
electron in two states at once. It may be the meaningful reading of the two forms of helium.

\emph{Helium.} Take the two forms of helium as these two modes. A radial self-consistent RealQM solve
computes, at one Hartree level, the two shell rates $E_1=-1.73,\ E_2=-0.15$~Ha (inner $\langle r\rangle=0.76$,
outer $4.9\,a_0$) and the total energies $E(1s2s)=-2.15$ and $E(1s^2)=-2.85$~Ha (exact $-2.9037$; the
residual is correlation, omitted at this level). Then
\begin{align}
\text{ortho, sloshing:}&\quad E(1s2s)-E(1s^2)=0.70\ \text{Ha}=19.0\ \text{eV}
  &&(\text{between two full states; obs.\ }19.8),\\
\text{para, breathing:}&\quad |E_1-E_2|=1.58\ \text{Ha}=43\ \text{eV}
  &&(\text{between the two shells of one state}),
\end{align}
cleanly separated by a factor of two, the breathing lying at the inner-shell scale (near the He$^+$
Lyman-$\alpha$, $40.8$~eV). One caveat of naming: standard spectroscopy ties ortho/para to spin, and there it
is \emph{para} --- the singlet $1s^2$ --- that owns the ground state, whereas the structural reading here ties
them to the oscillation mode. Which computed frequency falls on which observed series is what fixes the
labels, and $19.8$ and $43$~eV are far enough apart that the data separates them without ambiguity.

\section{Where the fast free boundary is needed}
\label{sec:fast}

Charge crossing between two bound domains --- breathing (\S\ref{sec:spectra}) carried to where the boundary
must move with the flow --- is the regime the quasi-static boundary cannot cover: when the partition itself
reshapes as fast as the phase turns, there is no separation of scales and no factorisation \eqref{eq:ansatz}.
There one must let the free boundary be updated on the fast scale, moving with the flow,
$V_n=\mathbf j_j\!\cdot\!\mathbf n/\rho_j$, so that each electron keeps its charge instant by instant. Three
conditions --- density continuity, Neumann, and this flow condition --- then bear on one moving interface,
and they conflict: Neumann gives $\mathbf j_j\!\cdot\!\mathbf n=0$, forcing $V_n=0$. The fast moving boundary
is thus a genuine open problem, sidestepped by the quasi-static scheme only because that scheme declines to
pose it.

It is needed exactly when charge is redistributed \emph{bound-to-bound} --- between two confined domains ---
in about one electronic period, so that $V_n$ is of order the phase rate. The applied cases are few and
sharp:
\begin{itemize}
\item \emph{Attosecond charge migration} --- the flagship. After a sudden ionisation a hole hops across a
molecule in femto- to attoseconds, driven by electron correlation before the nuclei move; the boundary
between the hole region and the rest must migrate on the electronic scale.
\item \emph{Sudden bond cleavage} --- where a bond breaks faster than nuclear relaxation and \emph{which}
fragment keeps the shared electrons is decided by how the shared boundary resolves; the fast boundary motion
is the product-selecting branching itself.
\item \emph{Fast ion--atom charge exchange} --- a keV projectile passing on the electronic timescale, an
electron transferring mid-collision as the interface sweeps across.
\item \emph{Nonadiabatic transitions and conical intersections} --- where the Born--Oppenheimer separation,
which is the quasi-static assumption, breaks down by construction.
\end{itemize}
Two look-alikes do \emph{not} need it. Ordinary thermal chemistry is slow: the nuclei set the pace and the
domains follow adiabatically, so the quasi-static boundary is enough. And ionisation, detachment,
high-harmonic rescattering, Auger and radioactive decay send charge to the \emph{continuum}, across an outer
absorbing boundary --- a single domain leaking --- not across the shared inter-domain boundary; these are
fast but they are the decay-type outer-boundary problem, not the fork. The overdetermination of the free
boundary is therefore real physics exactly once one enters attosecond bound-to-bound dynamics, and nowhere
before it.

\section{Conclusion}
\label{sec:concl}

The passage from real-valued to time-dependent RealQM is as simple as it looks: let the wave functions be
complex and evolve each domain by the ordinary complex Schr\"odinger equation. Because the free boundary and
the Coulomb coupling see only the phase-blind density, the static construction is untouched; complexification
merely returns the phase the real theory had suppressed. Taking the geometry quasi-static and the phase fast
gives a two-scale theory --- a slow real-valued geometry of moving domains carrying fast per-domain phase
clocks --- whose energy differences supply exactly the time-dependent content (currents,
spectra, coherence) that a static real density lacks. No free-boundary equation of motion and no
overdetermination arise, because the boundary is never asked to chase the fast scale. The real-valued theory
remains the flagship; the complex theory is its direct time-dependent extension by a single step.

\begin{thebibliography}{9}
\bibitem{RealQM} C.~Johnson, \emph{Quantum Mechanics as Multiphase 3D Continuum Mechanics}, manuscript,
2026; \url{https://claes542.github.io/RealMolecule/gallery.html}.
\bibitem{RealNucleus} C.~Johnson, \emph{RealNucleus: Nuclear Binding as Dual Coulomb Confinement, without a
Strong or Weak Force}, manuscript, 2026.
\bibitem{ComplexRealQM} C.~Johnson, \emph{Complex Time-Dependent RealQM: Currents, Radiation, and Magnetism
from the Real-Space Continuum}, manuscript, 2026.
\end{thebibliography}

\end{document}
